\chapter{Conclusion}
\label{chap:conclusion}

This work developed a hot-started nullspace active-set method for the quadratic master
problems of regularized column generation, for multi-block linear programs whose linking
constraints may be equalities or inequalities of either direction.

\section{Summary}

The starting point was a \textbf{unified formulation} (\cref{chap:formulation}): surplus
variables bring every linking row into equality form, and the sign conditions of the
inequality rows become ordinary dual constraints of the surplus columns. Proximal
regularization of the dual then leads to a convex quadratic master problem
(\cref{chap:regularization}). Because all linking rows are equalities, its auxiliary variable
can be eliminated for all rows, and the result has only the convexity constraints and
nonnegativity bounds. Two identities underlie everything that follows: the duals of the
linking rows are an explicit affine function of the primal iterate, and the gradient of the
quadratic objective consists of the reduced costs of the columns and the duals of the
inequality rows.

\medskip

For the active-set method (\cref{chap:qpas}), the constraint matrix of a working set admits a
\textbf{reordering into five sets}: representatives, working-set columns, nullspace columns,
active surpluses and free surpluses. The surpluses violate the column generation property of
the convexity matrix, but this only adds two groups of columns with identity structure. The
reordering yields an explicit nullspace basis of vectors $e_j - e_{\rep(j)}$ and unit vectors,
which is never stored (\cref{chap:nullspace}). The reduced KKT system \textbf{decouples}: the
surplus part of the direction is explicit, and only the \textbf{effective reduced Hessian}
$\Heff = (\Aeff Z_x)^\top(\Aeff Z_x)$ has to be factorized. Here $\Aeff$ holds the equality rows
and the inequality rows with active surplus. The column part of the method is thus the method
for equality rows, applied to a linking matrix whose rows change with the surplus working set.
The Lagrange multipliers are the reduced costs and the dual signs, so the optimality test of
the master problem and the pricing test are the same test.

\medskip

Every change of the working set is a row addition, a row deletion or a rank-one modification
of a fixed-size $LDL^\top$ factorization (\cref{chap:updates}). A new column or a new
stability center leaves the factorization unchanged, which gives a hot start of every master
problem. Structural zeros left by the updates are controlled by periodic refactorization in
increasing index order. Singular reduced Hessians are resolved by an exact zero-curvature
step, which removes the dependency and keeps every working-set variable at its bound, for
released columns and released surpluses alike (\cref{chap:singular}).

\medskip

On generated crew scheduling instances (\cref{chap:results}), the method reproduces the optimum
of the full linear program and certifies it by the Lagrangian bound. Its master problems cost
between $0.01$ and $7$\,ms on average. That is two to three orders of magnitude less than a
general-purpose QP solver, and fast enough to make regularized column generation much faster
than classical column generation with a warm-started LP master: 7.7\,s instead of 650\,s on
the largest instance on which both were run. From medium-sized instances on, the complete
column generation run is even faster than solving the full linear program directly.
Inequality rows are handled at no extra cost and lead to shorter runs.


\section{Outlook}

The experiments point to three directions for further work.
\begin{itemize}
    \item \textbf{Dense factorization for large nullspaces.} The factor of $\Heff$ turned out
          to be dense. For large $|\Nx|$, a dense $LDL^\top$ (or $QR$) factorization with
          rank-one updates would avoid the overhead of sparse data structures and would make
          the periodic refactorization unnecessary.
    \item \textbf{Adaptive stabilization in the proximal phase.} The long tail of center-only
          iterations is a property of the proximal point method with a fixed $\delta$
          (\cref{prop:center-only}). Since $\delta$ does not enter the factorization, it can be
          decreased in this phase at no cost to the linear algebra.
    \item \textbf{Range rows and inexact master solves.} Linking rows with two finite bounds
          need surpluses with upper bounds, which fit into the same framework as an additional
          working-set state. Moreover, since the multipliers are the reduced costs, pricing
          could be called at any subspace minimizer instead of only at the optimum of the
          master problem. This would turn column generation into a single active-set method
          over all columns.
\end{itemize}
